%======================================================================
%  GUÍA DE EJERCICIOS · EL SISTEMA SOLAR EN EL SIMULADOR
%  --------------------------------------------------------------------
%  Compilar con: pdflatex guia-sistema-solar.tex (dos veces)
%  Licencia: GNU GPL v3
%======================================================================
\documentclass[11pt,a4paper]{article}

%----------------------------------------------------------------------
%  PAQUETES
%----------------------------------------------------------------------
\usepackage[utf8]{inputenc}
\usepackage[T1]{fontenc}
\usepackage[spanish,es-tabla,es-noindentfirst]{babel}
\usepackage{lmodern}
\usepackage{microtype}
\usepackage[a4paper,margin=2.5cm]{geometry}
\usepackage{amsmath,amssymb,amsthm}
\usepackage{siunitx}
\usepackage{booktabs}
\usepackage{array}
\usepackage{tabularx}
\usepackage{longtable}
\usepackage{multirow}
\usepackage{graphicx}
\usepackage{xcolor}
\usepackage{hyperref}
\usepackage{enumitem}
\usepackage{fancyhdr}
\usepackage{titlesec}
\usepackage{tcolorbox}
\usepackage{setspace}
\usepackage[font=small,labelfont=bf]{caption}
\usepackage{lastpage}
\usepackage{eso-pic}
\usepackage{mdframed}

\tcbuselibrary{theorems,skins,breakable}

%----------------------------------------------------------------------
%  COLORES INSTITUCIONALES
%----------------------------------------------------------------------
\definecolor{azulU}{HTML}{1B3A6B}
\definecolor{azulClaro}{HTML}{4A6FA5}
\definecolor{grisFondo}{HTML}{F4F6F9}
\definecolor{grisBorde}{HTML}{C9D2DD}
\definecolor{naranja}{HTML}{E08B2D}

%----------------------------------------------------------------------
%  HIPERVÍNCULOS
%----------------------------------------------------------------------
\hypersetup{
  colorlinks=false,
  pdftitle={Estudio del Sistema Solar mediante el Simulador},
  pdfauthor={MTI. Juan Carlos Santiago Sánchez},
  pdfsubject={Astronomía · Mecánica Celeste · Física},
  pdfkeywords={Kepler, gravitación universal, Sistema Solar, simulador}
}

%----------------------------------------------------------------------
%  DATOS DEL AUTOR  ←←←  MODIFICA AQUÍ
%----------------------------------------------------------------------
\newcommand{\autor}{MTI. Juan Carlos Santiago Sánchez}
\newcommand{\correo}{juancarlos.santiago.cb161@dgeti.sems.gob.mx}
\newcommand{\fechaDoc}{19 de octubre de 2026}
\newcommand{\institucion}{DGETI -- SEMS}
\newcommand{\facultad}{\strut}
\newcommand{\asignatura}{Física y Astronomía}
\newcommand{\profesor}{MTI. Juan Carlos Santiago Sánchez}
\newcommand{\ciudad}{\strut}
\pagestyle{plain}
\titleformat{\section}{\Large\bfseries}{\thesection.}{0.6em}{}
\titleformat{\subsection}{\large\bfseries}{\thesubsection}{0.5em}{}
\titleformat{\subsubsection}{\normalsize\bfseries}{\thesubsubsection}{0.5em}{}

\newenvironment{cajaTeoria}{\begin{quote}\small}{\end{quote}}
\newenvironment{cajaImportante}{\begin{quote}\small}{\end{quote}}
\newenvironment{cajaEjercicio}[1][]{\par\medskip\noindent\textbf{#1}\par\smallskip}{\par\medskip}

%----------------------------------------------------------------------
%  ENTORNOS MATEMÁTICOS
%----------------------------------------------------------------------
\theoremstyle{definition}
\newtheorem{ley}{Ley}
\newtheorem{definicion}{Definición}
\newtheorem{ejemplo}{Ejemplo}

%----------------------------------------------------------------------
%  ESPACIADO
%----------------------------------------------------------------------
\setlength{\parskip}{0.35em}
\setlength{\parindent}{1.2em}
\setstretch{1.12}

%----------------------------------------------------------------------
%  COMANDOS ÚTILES
%----------------------------------------------------------------------
\newcommand{\UA}{\si{\astronomicalunit}}
\newcommand{\Msun}{\ensuremath{M_{\odot}}}
\newcommand{\Mearth}{\ensuremath{M_{\oplus}}}
\newcommand{\Rearth}{\ensuremath{R_{\oplus}}}
\newcommand{\GMsun}{\ensuremath{GM_{\odot}}}
\newcommand{\G}{\ensuremath{G}}
\newcommand{\vel}{\ensuremath{v}}
\newcommand{\velesc}{\ensuremath{v_{\text{esc}}}}

%======================================================================
\begin{document}

\begin{center}
{\LARGE\bfseries ESTUDIO DEL SISTEMA SOLAR Y SU MODELACIÓN ORBITAL}\\[0.4em]
{\large Guía didáctica: marco conceptual, ejercicios, aplicaciones y datos de referencia}\\[1.2em]
\autor\\
\correo\\[0.7em]
\fechaDoc
\end{center}

\tableofcontents
\newpage

%======================================================================
\section{Presentación}
%======================================================================

El \textbf{Simulador del Sistema Solar} (archivo \texttt{index.html}) es una aplicación autocontenida en HTML y JavaScript que permite explorar, en tiempo real y en un navegador, la estructura y dinámica del Sistema Solar. Para cada cuerpo --el Sol, los ocho planetas y dieciséis satélites principales-- el simulador muestra en un panel lateral:

\begin{itemize}[leftmargin=*,itemsep=2pt]
  \item masa y radio,
  \item semieje mayor, excentricidad e inclinación,
  \item período orbital y distancia al primario,
  \item velocidad orbital instantánea (ecuación vis-viva),
  \item aceleración gravitatoria.
\end{itemize}

\subsection*{Propósito general}

Comprender y aplicar las leyes de Kepler y la gravitación newtoniana para interpretar datos orbitales del Sistema Solar mediante un simulador, reconociendo los supuestos del modelo y contrastando los resultados con valores astronómicos de referencia.

\subsection*{Competencias esperadas}

Al completar esta guía, el estudiante será capaz de:

\begin{enumerate}[leftmargin=*,itemsep=2pt]
  \item \textbf{Registrar y organizar} datos astronómicos reales usando una herramienta digital.
  \item \textbf{Formular} las tres leyes de Kepler y reconocerlas en la dinámica observada.
  \item \textbf{Aplicar} la ley de gravitación universal para calcular fuerzas, aceleraciones y velocidades.
  \item \textbf{Verificar} la tercera ley de Kepler con datos reales y estimar errores.
  \item \textbf{Predecir} el comportamiento de cuerpos no presentes en el simulador mediante cálculo.
  \item \textbf{Extender} el catálogo del simulador añadiendo nuevos astros con datos coherentes.
\end{enumerate}

\subsection*{Cómo usar esta guía}

Cada bloque contiene una \textbf{referencia teórica breve} o remite al marco teórico general (Sección~\ref{sec:marco}). Los ejercicios están pensados para resolverse \textbf{con el simulador abierto en paralelo}. Las tablas vacías se completan leyendo el panel de datos del simulador. Los valores de referencia para el docente aparecen en el Apéndice~\ref{ap:datos}.

\subsection*{Constantes y convenciones}

\begin{table}[h]
\centering
\caption{Constantes físicas usadas en la guía.}
\label{tab:constantes}
\begin{tabular}{@{}lll@{}}
\toprule
\textbf{Constante} & \textbf{Símbolo} & \textbf{Valor} \\
\midrule
Unidad astronómica & \si{\astronomicalunit} & \SI{1.495978707e11}{\metre} \\
Masa de la Tierra & \Mearth & \SI{5.97237e24}{\kilo\gram} \\
Radio de la Tierra & \Rearth & \SI{6371}{\kilo\metre} \\
$\GMsun$ & \GMsun & \SI{1.32712440018e20}{\cubic\metre\per\square\second} \\
Año juliano & 1 año juliano & \SI{365.25}{\day} = \SI{3.15576e7}{\second} \\
Día & \SI{1}{\day} & \SI{86400}{\second} \\
Constante de gravitación & $\G$ & \SI{6.67430e-11}{\newton\square\metre\per\square\kilo\gram} \\
\bottomrule
\end{tabular}
\end{table}

\begin{cajaImportante}
\textbf{Nota sobre unidades.} En el simulador, los parámetros gravitacionales \texttt{GM} están en unidades del SI (\si{\cubic\metre\per\square\second}). El panel muestra distancias y semiejes en \si{\astronomicalunit} y en \si{\kilo\metre}; el estudiante deberá convertir al SI cuando lo requiera el cálculo.
\end{cajaImportante}

%======================================================================
\section{Marco teórico}\label{sec:marco}
%======================================================================

\subsection{Leyes de Kepler (1609--1619)}

Johannes Kepler dedujo empíricamente, a partir de las observaciones de Tycho Brahe, tres leyes que describen el movimiento planetario alrededor del Sol.

\begin{ley}[Primera ley: órbitas elípticas]
En el modelo de dos cuerpos, cada planeta describe una órbita elíptica y el centro de masa del sistema ocupa uno de sus focos. Cuando la masa del planeta es mucho menor que la solar, este foco se aproxima al centro del Sol.
\end{ley}

Una elipse queda definida por el \textbf{semieje mayor} $a$, el \textbf{semieje menor} $b = a\sqrt{1-e^{2}}$ y la \textbf{excentricidad} $e$. La distancia al Sol varía según la posición en la órbita:
\[
  r_p = a(1-e) \quad\text{(perihelio)}, \qquad
  r_a = a(1+e) \quad\text{(afelio)}.
\]

\begin{ley}[Segunda ley: ley de las áreas]
La línea que une al planeta con el Sol barre áreas iguales en tiempos iguales.
\end{ley}

Consecuencia directa: el planeta se mueve \textbf{más rápido en el perihelio} y \textbf{más lento en el afelio}. En términos modernos, esto equivale a la \textbf{conservación del momento angular} de cada planeta respecto al Sol.

\begin{ley}[Tercera ley: ley armónica]
El cuadrado del período orbital es proporcional al cubo del semieje mayor:
\[
  T^{2} = k\,a^{3}.
\]
Para el Sistema Solar, con $T$ en años y $a$ en \si{\astronomicalunit} alrededor del Sol, $k=1$:
\[
  T^{2} = a^{3}.
\]
\end{ley}

Para una órbita elíptica bajo la gravitación newtoniana, la forma general, considerando también la masa $m$ del cuerpo orbitante, es:
\[
  T^{2} = \frac{4\pi^{2}}{GM}\,a^{3}.
\]

\subsection{Ley de gravitación universal (Newton, 1687)}

\begin{ley}[Gravitación universal]
Dos cuerpos cualesquiera se atraen con una fuerza proporcional al producto de sus masas e inversamente proporcional al cuadrado de la distancia que los separa:
\[
  F = G\,\frac{M\,m}{r^{2}}.
\]
\end{ley}

Es habitual agrupar $\mu \equiv GM$ como el \textbf{parámetro gravitacional estándar} del cuerpo central. Para el Sol, $\GMsun = \SI{1.32712440018e20}{\cubic\metre\per\square\second}$.

\subsubsection*{Deducción de la tercera ley de Kepler}

Para una masa $m$ en órbita circular de radio $r$ alrededor de $M\gg m$, la fuerza gravitatoria es la fuerza centrípeta:
\[
  \frac{GMm}{r^{2}} = \frac{mv^{2}}{r}
  \;\Longrightarrow\;
  v^{2} = \frac{GM}{r}.
\]
Sustituyendo $v = 2\pi r/T$:
\[
  \left(\frac{2\pi r}{T}\right)^{2} = \frac{GM}{r}
  \;\Longrightarrow\;
  T^{2} = \frac{4\pi^{2}}{GM}\,r^{3},
\]
que coincide con la ley armónica para $k = 4\pi^{2}/(GM)$.

\subsection{Cantidades orbitales derivadas}

\subsubsection*{Velocidad orbital: ecuación vis-viva}

Para una órbita elíptica de semieje mayor $a$, la velocidad a distancia $r$ del cuerpo central es
\begin{equation}
  v = \sqrt{GM\left(\frac{2}{r} - \frac{1}{a}\right)}.
  \label{eq:visviva}
\end{equation}

Casos particulares:
\[
  v_{\text{circ}} = \sqrt{\frac{GM}{a}}, \qquad
  v_{p} = \sqrt{\frac{GM(1+e)}{a(1-e)}}, \qquad
  v_{a} = \sqrt{\frac{GM(1-e)}{a(1+e)}}.
\]

\subsubsection*{Aceleración gravitatoria}

\[
  g_{\text{orb}} = \frac{GM}{r^{2}}, \qquad
  g_{\text{sup}} = \frac{GM}{R^{2}}.
\]

\subsubsection*{Velocidad de escape}

\[
  \velesc = \sqrt{\frac{2GM}{R}}.
\]

\subsubsection*{Velocidad angular media}

\[
  \omega = \frac{2\pi}{T}.
\]

\subsection{Relación entre la teoría y el simulador}

\begin{table}[h]
\centering
\caption{Correspondencia entre conceptos teóricos y elementos del simulador.}
\label{tab:correspondencia}
\begin{tabular}{@{}p{0.45\linewidth}p{0.5\linewidth}@{}}
\toprule
\textbf{Concepto teórico} & \textbf{Ubicación en el simulador} \\
\midrule
Semieje mayor $a$ & Panel de datos -- \emph{Semieje mayor} \\
Excentricidad $e$ & Panel de datos -- \emph{Excentricidad} \\
Inclinación $i$ & Panel de datos -- \emph{Inclinación} \\
Período $T$ & Panel de datos -- \emph{Periodo orbital} \\
Distancia al primario $r$ & Panel de datos -- \emph{distancia al cuerpo primario} \\
Velocidad orbital $v$ & Panel de datos -- \emph{Velocidad orbital} \\
Aceleración gravitatoria & Panel de datos -- \emph{Aceleración grav.} \\
Masa $M$ y radio $R$ & Panel de datos -- \emph{Masa}, \emph{Radio} \\
Parámetro $\mu = GM$ & Código fuente (\texttt{CATALOG}) \\
Forma elíptica de la órbita & Botón \emph{Órbitas} activado \\
\bottomrule
\end{tabular}
\end{table}

\begin{cajaImportante}
\textbf{Sugerencia didáctica.} Activa los \emph{rastros} y observa cómo la velocidad de un planeta aumenta cerca del perihelio. Es la \textbf{segunda ley de Kepler} en acción.
\end{cajaImportante}

%======================================================================
\section{Bloque 0: Familiarización con el simulador}
%======================================================================

\begin{cajaEjercicio}[Ejercicio 0.1 -- Primer contacto]
Abre el simulador y responde:
\begin{enumerate}[label=\alph*),leftmargin=*]
  \item ¿Cuántos planetas y cuántos satélites puedes ver en total?
  \item ¿Qué cuerpo tiene la órbita de mayor radio? ¿Y la de menor?
  \item Activa la \emph{inclinación de vista} a $60^{\circ}$. ¿Qué planetas se ven claramente fuera del plano medio?
  \item Selecciona el \textbf{Sol}. Anota su radio, masa y \texttt{GM}.
\end{enumerate}
\end{cajaEjercicio}

\begin{cajaEjercicio}[Ejercicio 0.2 -- Lectura del panel de datos]
Selecciona la \textbf{Tierra} y completa la siguiente tabla:

\begin{center}
\begin{tabular}{@{}lc@{}}
\toprule
\textbf{Dato} & \textbf{Valor} \\
\midrule
Masa & \hspace{4cm} \\
Radio & \\
Semieje mayor $a$ & \\
Excentricidad $e$ & \\
Inclinación $i$ & \\
Período orbital $T$ & \\
Distancia al Sol $r$ & \\
Velocidad orbital $v$ & \\
Aceleración gravitatoria $g$ & \\
\bottomrule
\end{tabular}
\end{center}
\end{cajaEjercicio}

\begin{cajaEjercicio}[Ejercicio 0.3 -- La escala]
Observa el HUD superior izquierdo:
\begin{enumerate}[label=\alph*),leftmargin=*]
  \item ¿Cuántos \si{\kilo\metre} representa un píxel en la vista \emph{Sistema completo}?
  \item Haz zoom hasta ver la órbita de la Luna. ¿Cuántos \si{\kilo\metre} representa un píxel ahora?
  \item ¿Por qué el simulador exagera el tamaño de los planetas a escala de sistema completo?
\end{enumerate}
\end{cajaEjercicio}

%======================================================================
\section{Bloque 1: Tablas de datos -- Planetas}
%======================================================================

\begin{cajaEjercicio}[Ejercicio 1.1 -- Tabla comparativa de los ocho planetas]
Completa la tabla leyendo el panel de datos del simulador. Para las columnas \emph{Perihelio} y \emph{Afelio}, calcúlalas con $r_p = a(1-e)$ y $r_a = a(1+e)$.

\begin{longtable}{@{}l*{9}{c}@{}}
\caption{Tabla comparativa de los planetas.} \\
\toprule
\textbf{Planeta} & $R$\,(\si{\kilo\metre}) & $M$\,(\si{\kilo\gram}) & $a$\,(\si{\astronomicalunit}) & $e$ & $i$\,($^{\circ}$) & $T$\,(\si{\day}) & $r_p$ & $r_a$ & $v$\,(\si{\kilo\metre\per\second}) \\
\midrule
\endfirsthead
\toprule
\textbf{Planeta} & $R$ & $M$ & $a$ & $e$ & $i$ & $T$ & $r_p$ & $r_a$ & $v$ \\
\midrule
\endhead
Mercurio & & & & & & & & & \\
Venus    & & & & & & & & & \\
Tierra   & & & & & & & & & \\
Marte    & & & & & & & & & \\
Júpiter  & & & & & & & & & \\
Saturno  & & & & & & & & & \\
Urano    & & & & & & & & & \\
Neptuno  & & & & & & & & & \\
\bottomrule
\end{longtable}
\end{cajaEjercicio}

\textbf{Cuestiones:}

\begin{enumerate}[label=\alph*),leftmargin=*]
  \item ¿Qué planeta tiene la órbita más excéntrica? ¿Y la menos excéntrica?
  \item ¿Qué planeta tarda más en completar una órbita? ¿Por qué?
  \item Ordena los planetas por velocidad orbital de mayor a menor. ¿Qué patrón observas?
  \item ¿Cuántas veces es más masivo Júpiter que la Tierra? ¿Y cuántas veces cabe la Tierra en el radio de Júpiter?
\end{enumerate}

\begin{cajaEjercicio}[Ejercicio 1.2 -- Clasificación y densidad]
A partir de los datos de la tabla anterior:
\begin{enumerate}[label=\alph*),leftmargin=*]
  \item Clasifica los planetas en \emph{rocosos} y \emph{gigantes gaseosos}.
  \item Calcula la densidad media de cada planeta con
  \[
    \rho = \frac{M}{\tfrac{4}{3}\pi R^{3}} \quad [\si{\kilo\gram\per\cubic\metre}].
  \]
  \item Compara con la densidad del agua ($\SI{1000}{\kilo\gram\per\cubic\metre}$) y con la de la Tierra ($\approx \SI{5500}{\kilo\gram\per\cubic\metre}$).
\end{enumerate}
\end{cajaEjercicio}

\begin{cajaEjercicio}[Ejercicio 1.3 -- Mercurio frente a Neptuno]
Completa la tabla y responde:

\begin{center}
\begin{tabular}{@{}lcc@{}}
\toprule
 & \textbf{Mercurio} & \textbf{Neptuno} \\
\midrule
$a$ (\si{\astronomicalunit}) & & \\
$T$ (años) & & \\
$v$ orbital media (\si{\kilo\metre\per\second}) & & \\
\bottomrule
\end{tabular}
\end{center}

\begin{enumerate}[label=\alph*),leftmargin=*]
  \item ¿Cuántas veces es mayor el semieje de Neptuno que el de Mercurio?
  \item ¿Cuántas veces es mayor el período de Neptuno que el de Mercurio?
  \item ¿Se cumple aproximadamente $T^{2}\propto a^{3}$?
\end{enumerate}
\end{cajaEjercicio}

%======================================================================
\section{Bloque 2: Tablas de datos -- Satélites}
%======================================================================

\begin{cajaEjercicio}[Ejercicio 2.1 -- Lunas principales]
Completa la tabla leyendo el panel del simulador. Observa que, para las lunas, $a$ se mide \textbf{respecto de su planeta}, no del Sol.

\begin{longtable}{@{}llcccccc@{}}
\caption{Lunas principales del Sistema Solar.} \\
\toprule
\textbf{Satélite} & \textbf{Padre} & $R$\,(\si{\kilo\metre}) & $M$\,(\si{\kilo\gram}) & $a$\,(\si{\kilo\metre}) & $e$ & $i$\,($^{\circ}$) & $T$\,(\si{\day}) \\
\midrule
\endfirsthead
\toprule
\textbf{Satélite} & \textbf{Padre} & $R$ & $M$ & $a$ & $e$ & $i$ & $T$ \\
\midrule
\endhead
Luna      & Tierra   & & & & & & \\
Fobos     & Marte    & & & & & & \\
Deimos    & Marte    & & & & & & \\
Ío        & Júpiter  & & & & & & \\
Europa    & Júpiter  & & & & & & \\
Ganímedes & Júpiter  & & & & & & \\
Calisto   & Júpiter  & & & & & & \\
Encélado  & Saturno  & & & & & & \\
Rea       & Saturno  & & & & & & \\
Titán     & Saturno  & & & & & & \\
Jápeto    & Saturno  & & & & & & \\
Ariel     & Urano    & & & & & & \\
Umbriel   & Urano    & & & & & & \\
Titania   & Urano    & & & & & & \\
Oberón    & Urano    & & & & & & \\
Tritón    & Neptuno  & & & & & & \\
\bottomrule
\end{longtable}
\end{cajaEjercicio}

\textbf{Cuestiones:}

\begin{enumerate}[label=\alph*),leftmargin=*]
  \item ¿Qué satélite tiene la órbita más inclinada respecto al ecuador de su planeta?
  \item Ganímedes es el satélite más grande del Sistema Solar. ¿Cómo se compara su radio con el de Mercurio?
  \item ¿Por qué las lunas de Júpiter completan sus órbitas en pocos días, mientras que las de Urano tardan semanas?
\end{enumerate}

\begin{cajaEjercicio}[Ejercicio 2.2 -- Satélites y planetas]
\begin{enumerate}[label=\alph*),leftmargin=*]
  \item Compara la Luna y Titán con los planetas interiores en tamaño y masa.
  \item ¿Qué satélite tarda más en orbitar su planeta? ¿Por qué?
  \item La Luna se aleja de la Tierra a razón de $3.8\,\mathrm{cm\,a^{-1}}$. Si $a_{\text{Luna}} = \SI{384400}{\kilo\metre}$, ¿en cuántos años se duplicaría esa distancia, suponiendo tasa constante? ¿Es un tiempo comparable a la edad del Universo ($\SI{1.38e10}{}$ años)?
\end{enumerate}
\end{cajaEjercicio}

%======================================================================
\section{Bloque 3: Las leyes de Kepler}
%======================================================================

\begin{cajaEjercicio}[Ejercicio 3.1 -- Primera ley en el simulador]
Con la vista inclinada a $0^{\circ}$ y zoom amplio sobre un planeta interior:
\begin{enumerate}[label=\alph*),leftmargin=*]
  \item Observa la órbita de Mercurio. ¿Se ve claramente elíptica o casi circular? ¿Por qué?
  \item Repite con Marte. ¿Notas más o menos excentricidad? Justifica con los valores de $e$.
  \item Selecciona Mercurio y observa su distancia al Sol durante un tiempo. Anota los valores mínimo y máximo.
\end{enumerate}
\end{cajaEjercicio}

\begin{cajaEjercicio}[Ejercicio 3.2 -- Segunda ley: velocidad y distancia]
\begin{enumerate}[label=\alph*),leftmargin=*]
  \item Selecciona Mercurio y observa cómo cambia su velocidad orbital.
  \item Anota la velocidad cerca del perihelio y del afelio.
  \item Calcula la relación $v_{p}/v_{a}$ y compárala con el valor teórico:
  \[
    \frac{v_{p}}{v_{a}} = \frac{1+e}{1-e}.
  \]
  \item ¿Se cumple la ley de las áreas? Explica por qué esa relación se deduce de la conservación del momento angular.
\end{enumerate}
\end{cajaEjercicio}

\begin{cajaEjercicio}[Ejercicio 3.3 -- Tercera ley: la tabla armónica]
Usando los datos del Ejercicio 1.1, completa:

\begin{longtable}{@{}lccccc@{}}
\caption{Verificación de la tercera ley de Kepler.} \\
\toprule
\textbf{Planeta} & $a$\,(\si{\astronomicalunit}) & $T$\,(años) & $a^{3}$ & $T^{2}$ & $T^{2}/a^{3}$ \\
\midrule
\endfirsthead
\toprule
\textbf{Planeta} & $a$ & $T$ & $a^{3}$ & $T^{2}$ & $T^{2}/a^{3}$ \\
\midrule
\endhead
Mercurio & & & & & \\
Venus    & & & & & \\
Tierra   & & & & & \\
Marte    & & & & & \\
Júpiter  & & & & & \\
Saturno  & & & & & \\
Urano    & & & & & \\
Neptuno  & & & & & \\
\bottomrule
\end{longtable}
\end{cajaEjercicio}

\textbf{Preguntas:}
\begin{enumerate}[label=\alph*),leftmargin=*]
  \item ¿Todos los cocientes $T^{2}/a^{3}$ son iguales? ¿A qué valor se aproximan?
  \item ¿Qué unidades se usaron para $T$ y $a$? ¿Por qué importa?
  \item Si $T$ se expresara en días en lugar de años, ¿cuál sería el valor de la constante?
\end{enumerate}

\begin{cajaEjercicio}[Ejercicio 3.4 -- Extensión: ¿se cumple para las lunas?]
Para el sistema \textbf{Júpiter--lunas}, aplica la tercera ley con $k = 4\pi^{2}/(GM_{J})$:
\begin{enumerate}[label=\alph*),leftmargin=*]
  \item Con $GM_{J} = \SI{1.26686534e17}{\cubic\metre\per\square\second}$ y los datos de las lunas galileanas (Ío, Europa, Ganímedes, Calisto), calcula $T^{2}/a^{3}$ en unidades del SI.
  \item Compara con $4\pi^{2}/(GM_{J})$. ¿Se cumple?
  \item ¿Por qué la constante para Júpiter es distinta que para el Sol?
\end{enumerate}
\end{cajaEjercicio}

\begin{cajaEjercicio}[Ejercicio 3.5 -- Predicción]
Supón que descubres un nuevo planeta en el Sistema Solar con $a=\SI{4.0}{\astronomicalunit}$. Sin usar el simulador:
\begin{enumerate}[label=\alph*),leftmargin=*]
  \item ¿Cuál sería su período orbital en años?
  \item ¿Cuál sería su velocidad orbital media aproximada, en \si{\kilo\metre\per\second}?
  \item Verifica tus resultados comparando con planetas de $a$ similares (Marte: $\SI{1.52}{\astronomicalunit}$; Júpiter: $\SI{5.20}{\astronomicalunit}$).
\end{enumerate}
\end{cajaEjercicio}

%======================================================================
\section{Bloque 4: Gravitación universal}
%======================================================================

\begin{cajaEjercicio}[Ejercicio 4.1 -- Fuerza gravitatoria]
Calcula la fuerza que ejerce el Sol sobre cada planeta con
\[
  F = G\,\frac{\Msun\, m}{r^{2}}, \qquad
  \Msun = \SI{1.9885e30}{\kilo\gram}.
\]

\begin{center}
\begin{tabular}{@{}lccc@{}}
\toprule
\textbf{Planeta} & $m$\,(\si{\kilo\gram}) & $a$\,(\si{\metre}) & $F$\,(\si{\newton}) \\
\midrule
Mercurio & & & \\
Tierra   & & & \\
Júpiter  & & & \\
Neptuno  & & & \\
\bottomrule
\end{tabular}
\end{center}

\textbf{Preguntas:}
\begin{enumerate}[label=\alph*),leftmargin=*]
  \item ¿Sobre qué planeta actúa la fuerza gravitatoria más intensa? ¿Y la más débil?
  \item ¿Cómo influyen simultáneamente la masa del planeta y su distancia al Sol en la fuerza? Compara el resultado para Neptuno y Mercurio.
  \item ¿Qué fuerza ejerce cada planeta sobre el Sol? ¿Son iguales o distintas? (\emph{Tercera ley de Newton}).
\end{enumerate}
\end{cajaEjercicio}

\begin{cajaEjercicio}[Ejercicio 4.2 -- Aceleración gravitatoria]
Calcula la aceleración que el Sol produce sobre cada planeta a la distancia media:
\[
  g = \frac{\GMsun}{a^{2}}.
\]
\begin{center}
\begin{tabular}{@{}lcc@{}}
\toprule
\textbf{Planeta} & $a$\,(\si{\metre}) & $g$\,(\si{\metre\per\square\second}) \\
\midrule
Mercurio & & \\
Tierra   & & \\
Júpiter  & & \\
Neptuno  & & \\
\bottomrule
\end{tabular}
\end{center}

\begin{enumerate}[label=\alph*),leftmargin=*]
  \item Compara con la \emph{aceleración gravitatoria} que muestra el simulador. ¿Coinciden?
  \item ¿Por qué la aceleración es mayor para Mercurio que para Neptuno?
  \item ¿Qué relación hay entre esta aceleración y la velocidad orbital?
\end{enumerate}
\end{cajaEjercicio}

\begin{cajaEjercicio}[Ejercicio 4.3 -- Gravedad superficial]
Calcula la gravedad superficial $g_{\text{sup}} = GM/R^{2}$ de los siguientes cuerpos:

\begin{center}
\begin{tabular}{@{}lcccc@{}}
\toprule
\textbf{Cuerpo} & $GM$\,(\si{\cubic\metre\per\square\second}) & $R$\,(\si{\metre}) & $g_{\text{sup}}$\,(\si{\metre\per\square\second}) & \textbf{vs.\ Tierra} \\
\midrule
Luna    & & & & \\
Marte   & & & & \\
Júpiter & & & & \\
Titán   & & & & \\
Ío      & & & & \\
\bottomrule
\end{tabular}
\end{center}

\begin{enumerate}[label=\alph*),leftmargin=*]
  \item ¿En qué cuerpo un astronauta de \SI{70}{\kilo\gram} pesaría menos? ¿En cuál más?
  \item ¿Por qué la gravedad superficial de Júpiter (a pesar de su enorme masa) no es proporcionalmente gigantesca?
\end{enumerate}
\end{cajaEjercicio}

\begin{cajaEjercicio}[Ejercicio 4.4 -- Velocidad de escape]
Calcula $\velesc = \sqrt{2GM/R}$ para los siguientes cuerpos:

\begin{center}
\begin{tabular}{@{}lccc@{}}
\toprule
\textbf{Cuerpo} & $GM$\,(\si{\cubic\metre\per\square\second}) & $R$\,(\si{\metre}) & $\velesc$\,(\si{\kilo\metre\per\second}) \\
\midrule
Tierra  & $3.986\times10^{14}$ & $6.371\times10^{6}$ & \\
Luna    & $4.905\times10^{12}$ & $1.737\times10^{6}$ & \\
Marte   & $4.283\times10^{13}$ & $3.390\times10^{6}$ & \\
Júpiter & $1.267\times10^{17}$ & $6.991\times10^{7}$ & \\
Sol     & $1.327\times10^{20}$ & $6.963\times10^{8}$ & \\
\bottomrule
\end{tabular}
\end{center}

\begin{enumerate}[label=\alph*),leftmargin=*]
  \item ¿Por qué la Luna no tiene atmósfera significativa?
  \item ¿Qué implicaría intentar despegar del Sol? ¿Es viable con la tecnología actual?
\end{enumerate}
\end{cajaEjercicio}

\begin{cajaEjercicio}[Ejercicio 4.5 -- Peso aparente y órbitas]
\begin{enumerate}[label=\alph*),leftmargin=*]
  \item ¿Por qué los astronautas en la Estación Espacial Internacional (a $\sim\SI{400}{\kilo\metre}$ de altura) flotan, si la gravedad terrestre allí es casi la misma que en superficie?
  \item Calcula $g$ a \SI{400}{\kilo\metre} de altura sobre la Tierra. Compáralo con $g$ en superficie.
\end{enumerate}
\end{cajaEjercicio}

%======================================================================
\section{Bloque 5: Cantidades orbitales y aplicaciones}
%======================================================================

\begin{cajaEjercicio}[Ejercicio 5.1 -- Perihelio, afelio y velocidad extrema]
Para Mercurio ($a = \SI{0.387098}{\astronomicalunit}$, $e=0.2056$, $\GMsun = \SI{1.327e20}{\cubic\metre\per\square\second}$):
\begin{enumerate}[label=\alph*),leftmargin=*]
  \item Calcula $r_{p}$ y $r_{a}$ en \si{\astronomicalunit} y en \si{\metre}.
  \item Calcula $v_{p}$ y $v_{a}$ con la ecuación vis-viva \eqref{eq:visviva}.
  \item Calcula la relación $v_{p}/v_{a}$ y compárala con $(1+e)/(1-e)$.
  \item Verifica en el simulador seleccionando Mercurio y observando su velocidad.
\end{enumerate}
\end{cajaEjercicio}

\begin{cajaEjercicio}[Ejercicio 5.2 -- Comparación de velocidades]
Calcula la velocidad orbital media (circular) de los ocho planetas:
\[
  v_{\text{circ}} = \sqrt{\GMsun/a}.
\]
\begin{center}
\begin{tabular}{@{}lcc@{}}
\toprule
\textbf{Planeta} & $a$\,(\si{\metre}) & $v_{\text{circ}}$\,(\si{\kilo\metre\per\second}) \\
\midrule
Mercurio & & \\
Venus    & & \\
Tierra   & & \\
Marte    & & \\
Júpiter  & & \\
Saturno  & & \\
Urano    & & \\
Neptuno  & & \\
\bottomrule
\end{tabular}
\end{center}

\begin{enumerate}[label=\alph*),leftmargin=*]
  \item Compara con los valores del simulador. ¿Son cercanos? ¿Por qué no idénticos?
  \item Grafica $v$ en función de $a$. ¿Qué tipo de curva aparece?
  \item ¿Qué planeta se mueve más rápido? ¿Y más lento?
\end{enumerate}
\end{cajaEjercicio}

\begin{cajaEjercicio}[Ejercicio 5.3 -- Velocidad angular y ``años'']
\begin{enumerate}[label=\alph*),leftmargin=*]
  \item Calcula $\omega = 2\pi/T$ de cada planeta en \si{\radian por año}.
  \item ¿Cuántos ``años locales'' han transcurrido en los últimos \SI{100}{} años terrestres?
  \item Si vivieras en Marte \SI{1}{} año marciano, ¿cuántos años terrestres habrían pasado?
\end{enumerate}
\end{cajaEjercicio}

\begin{cajaEjercicio}[Ejercicio 5.4 -- Un satélite artificial]
Un satélite orbita la Tierra en una órbita circular a \SI{400}{\kilo\metre} de altura ($\Rearth = \SI{6371}{\kilo\metre}$, $GM_{\oplus} = \SI{3.986e14}{\cubic\metre\per\square\second}$):
\begin{enumerate}[label=\alph*),leftmargin=*]
  \item Calcula su velocidad orbital.
  \item Calcula su período orbital en minutos.
  \item ¿Cuántas órbitas completa en un día?
  \item Compara con la Luna (27.3 días). ¿Por qué tanta diferencia?
\end{enumerate}
\end{cajaEjercicio}

%======================================================================
\section{Bloque 6: Síntesis y proyecto final}
%======================================================================

\begin{cajaEjercicio}[Ejercicio 6.1 -- Diseña tu propio astro]
Vas a añadir un cuerpo al \texttt{CATALOG} del simulador. Debe ser un \textbf{planeta enano del cinturón de Kuiper} con $a\approx\SI{45}{\astronomicalunit}$, $e\approx 0.15$, $i\approx 20^{\circ}$.
\begin{enumerate}[label=\alph*),leftmargin=*]
  \item Estima su período orbital usando $T^{2}=a^{3}$.
  \item Calcula su velocidad orbital media en el afelio y en el perihelio.
  \item Propón valores de \texttt{radius}, \texttt{mass} y \texttt{GM} razonables para un cuerpo helado de $\sim\SI{1000}{\kilo\metre}$ de radio (inspírate en Plutón: $R=\SI{1188}{\kilo\metre}$, $M=\SI{1.303e22}{\kilo\gram}$).
  \item Escribe el objeto completo para añadirlo al \texttt{CATALOG} y verifícalo en el simulador.
\end{enumerate}
\end{cajaEjercicio}

\begin{cajaEjercicio}[Ejercicio 6.2 -- Investigación corta]
Elige un satélite del simulador y elabora un informe (1 página) que incluya: descubrimiento, características físicas, órbita, curiosidades y comparación con la Luna terrestre.
\end{cajaEjercicio}

\begin{cajaEjercicio}[Ejercicio 6.3 -- Pregunta abierta]
\emph{``Si el Sol tuviera el doble de masa, ¿cómo cambiarían los períodos orbitales, suponiendo distancias iguales?''}

Argumenta usando la tercera ley generalizada $T^{2} = 4\pi^{2}a^{3}/(GM)$. Calcula el factor de cambio y compruébalo con el simulador modificando temporalmente \texttt{GM\_SUN} en el código.
\end{cajaEjercicio}

%======================================================================
\appendix
\section{Constantes y fórmulas}\label{ap:constantes}
%======================================================================

\begin{cajaTeoria}
\textbf{Constantes fundamentales:}
\begin{align*}
  \SI{1}{\astronomicalunit} &= \SI{1.495978707e11}{\metre} \\
  1 año juliano             &= \SI{3.15576e7}{\second} \\
  \SI{1}{\day}              &= \SI{86400}{\second} \\
  G                         &= \SI{6.67430e-11}{\newton\square\metre\per\square\kilo\gram} \\
  \Msun                     &= \SI{1.9885e30}{\kilo\gram} \\
  \GMsun                    &= \SI{1.32712440018e20}{\cubic\metre\per\square\second} \\
  \Mearth                   &= \SI{5.97237e24}{\kilo\gram} \\
  \Rearth                   &= \SI{6.371e6}{\metre}
\end{align*}

\textbf{Fórmulas clave:}
\begin{align*}
  F &= \frac{GMm}{r^{2}} & &\text{(gravitación universal)} \\
  g &= \frac{GM}{r^{2}} & &\text{(aceleración gravitatoria)} \\
  v_{\text{circ}} &= \sqrt{\frac{GM}{r}} & &\text{(velocidad circular)} \\
  v &= \sqrt{GM\left(\frac{2}{r} - \frac{1}{a}\right)} & &\text{(vis-viva)} \\
  \velesc &= \sqrt{\frac{2GM}{R}} & &\text{(velocidad de escape)} \\
  T^{2} &= \frac{4\pi^{2}}{GM}\,a^{3} & &\text{(Kepler, generalizada)} \\
  r_{p} &= a(1-e), \quad r_{a} = a(1+e) & & \\
  \omega &= \frac{2\pi}{T} & & \\
  \rho &= \frac{M}{\tfrac{4}{3}\pi R^{3}} & &
\end{align*}
\end{cajaTeoria}

%======================================================================
\section{Datos de referencia para el docente}\label{ap:datos}
%======================================================================

\begin{cajaImportante}
Esta sección \textbf{no debe entregarse al estudiante} antes de que complete sus propias mediciones. Se incluye como clave de corrección y para agilizar la revisión de los ejercicios de cálculo.
\end{cajaImportante}

\subsection*{Planetas}

\begin{longtable}{@{}lccccccc@{}}
\caption{Datos de referencia de los planetas.} \\
\toprule
\textbf{Planeta} & $R$ (\si{\kilo\metre}) & $M$ (\si{\kilo\gram}) & $a$ (\si{\astronomicalunit}) & $e$ & $i$ ($^{\circ}$) & $T$ (\si{\day}) & $v_{\text{circ}}$ (\si{\kilo\metre\per\second}) \\
\midrule
\endfirsthead
\toprule
\textbf{Planeta} & $R$ & $M$ & $a$ & $e$ & $i$ & $T$ & $v_{\text{circ}}$ \\
\midrule
\endhead
Mercurio & 2439.7 & $3.3011\times10^{23}$ & 0.38710 & 0.2056 & 7.006 & 87.97 & 47.9 \\
Venus    & 6051.8 & $4.8675\times10^{24}$ & 0.72332 & 0.00676 & 3.398 & 224.70 & 35.0 \\
Tierra   & 6371.0 & $5.9724\times10^{24}$ & 1.00000 & 0.01673 & 0.001 & 365.26 & 29.8 \\
Marte    & 3389.5 & $6.4171\times10^{23}$ & 1.52371 & 0.09337 & 1.852 & 686.98 & 24.1 \\
Júpiter  & 69911  & $1.8982\times10^{27}$ & 5.20289 & 0.04839 & 1.304 & 4332.59 & 13.1 \\
Saturno  & 58232  & $5.6834\times10^{26}$ & 9.53668 & 0.05386 & 2.486 & 10759.2 & 9.6 \\
Urano    & 25362  & $8.6810\times10^{25}$ & 19.1892 & 0.04726 & 0.773 & 30685.4 & 6.8 \\
Neptuno  & 24622  & $1.0241\times10^{26}$ & 30.0699 & 0.00859 & 1.770 & 60189.0 & 5.4 \\
\bottomrule
\end{longtable}

\subsection*{Satélites principales}

\begin{longtable}{@{}llcccccc@{}}
\caption{Datos de referencia de los satélites.} \\
\toprule
\textbf{Satélite} & \textbf{Padre} & $R$ (\si{\kilo\metre}) & $M$ (\si{\kilo\gram}) & $a$ (\si{\kilo\metre}) & $e$ & $i$ ($^{\circ}$) & $T$ (\si{\day}) \\
\midrule
\endfirsthead
\toprule
\textbf{Satélite} & \textbf{Padre} & $R$ & $M$ & $a$ & $e$ & $i$ & $T$ \\
\midrule
\endhead
Luna      & Tierra  & 1737.4 & $7.342\times10^{22}$ & 384\,400   & 0.0549  & 5.145   & 27.32 \\
Fobos     & Marte   & 11.27  & $1.0659\times10^{16}$ & 9\,376    & 0.0000  & 1.08    & 0.319 \\
Deimos    & Marte   & 6.2    & $1.4762\times10^{15}$ & 23\,463   & 0.0000  & 1.79    & 1.26 \\
Ío        & Júpiter & 1821.6 & $8.932\times10^{22}$ & 421\,700   & 0.0041  & 0.05    & 1.77 \\
Europa    & Júpiter & 1560.8 & $4.800\times10^{22}$ & 671\,034   & 0.0094  & 0.47    & 3.55 \\
Ganímedes & Júpiter & 2634.1 & $1.4819\times10^{23}$ & 1\,070\,412 & 0.0013  & 0.20    & 7.15 \\
Calisto   & Júpiter & 2410.3 & $1.0759\times10^{23}$ & 1\,882\,709 & 0.0074  & 0.19    & 16.69 \\
Encélado  & Saturno & 252.1  & $1.080\times10^{20}$ & 237\,948   & 0.0047  & 0.02    & 1.37 \\
Rea       & Saturno & 763.8  & $2.306\times10^{21}$ & 527\,108   & 0.0013  & 0.35    & 4.52 \\
Titán     & Saturno & 2574.7 & $1.3452\times10^{23}$ & 1\,221\,870 & 0.0288  & 0.35    & 15.95 \\
Jápeto    & Saturno & 734.5  & $1.805\times10^{21}$ & 3\,560\,820 & 0.0286  & 15.47   & 79.32 \\
Ariel     & Urano   & 578.9  & $1.251\times10^{21}$ & 190\,900   & 0.0012  & 0.26    & 2.52 \\
Umbriel   & Urano   & 584.7  & $1.275\times10^{21}$ & 266\,000   & 0.0039  & 0.13    & 4.14 \\
Titania   & Urano   & 788.4  & $3.400\times10^{21}$ & 436\,300   & 0.0011  & 0.34    & 8.71 \\
Oberón    & Urano   & 761.4  & $3.076\times10^{21}$ & 583\,500   & 0.0014  & 0.06    & 13.46 \\
Tritón    & Neptuno & 1353.4 & $2.139\times10^{22}$ & 354\,759   & 0.000016 & 156.89 & 5.88 \\
\bottomrule
\end{longtable}

\subsection*{Verificación de la tercera ley de Kepler}

\begin{center}
\begin{tabular}{@{}lcccccc@{}}
\toprule
\textbf{Planeta} & $a$ (\si{\astronomicalunit}) & $T$ (años) & $a^{3}$ & $T^{2}$ & $T^{2}/a^{3}$ & \textbf{Error} \\
\midrule
Mercurio & 0.3871 & 0.2408 & 0.0580 & 0.0580 & 1.000 & $<0.1\%$ \\
Venus    & 0.7233 & 0.6152 & 0.3783 & 0.3785 & 1.000 & $<0.1\%$ \\
Tierra   & 1.0000 & 1.0000 & 1.0000 & 1.0000 & 1.000 & $<0.1\%$ \\
Marte    & 1.5237 & 1.8808 & 3.5378 & 3.5375 & 1.000 & $<0.1\%$ \\
Júpiter  & 5.2029 & 11.862 & 140.84 & 140.71 & 0.999 & $0.1\%$ \\
Saturno  & 9.5367 & 29.457 & 867.4  & 867.7  & 1.000 & $<0.1\%$ \\
Urano    & 19.189 & 84.011 & 7064   & 7058   & 0.999 & $0.1\%$ \\
Neptuno  & 30.070 & 164.79 & 27192  & 27155  & 0.999 & $0.1\%$ \\
\bottomrule
\end{tabular}
\end{center}

\noindent Todos los cocientes son $\approx 1$ con un error menor al $0.5\%$. La tercera ley de Kepler se cumple con los datos del simulador.

%======================================================================
\section{Rúbrica de evaluación}\label{ap:rubrica}
%======================================================================

\begin{longtable}{@{}p{0.22\linewidth}*{4}{p{0.17\linewidth}}@{}}
\caption{Rúbrica de evaluación sugerida.} \\
\toprule
\textbf{Criterio} & \textbf{Insuficiente (1)} & \textbf{Suficiente (2)} & \textbf{Bien (3)} & \textbf{Excelente (4)} \\
\midrule
\endfirsthead
\toprule
\textbf{Criterio} & \textbf{Insuficiente (1)} & \textbf{Suficiente (2)} & \textbf{Bien (3)} & \textbf{Excelente (4)} \\
\midrule
\endhead
Registro de datos & No completa tablas & Completa con errores & Completa correctamente & Completa y verifica coherencia \\
Aplicación de leyes de Kepler & No las identifica & Las identifica pero no aplica & Aplica con errores menores & Aplica correctamente y justifica \\
Cálculos de gravitación & No usa fórmulas & Usa fórmulas con errores & Cálculos correctos con unidades & Cálculos correctos, unidades y análisis \\
Unidades y notación científica & Confunde unidades & Algunas conversiones correctas & Conversiones correctas & Dominio completo del SI \\
Análisis y conclusiones & Sin conclusiones & Conclusiones superficiales & Conclusiones correctas & Conclusiones razonadas y conectadas con la teoría \\
Presentación y orden & Desordenado & Legible & Claro y ordenado & Claro, ordenado y con esquemas \\
\bottomrule
\end{longtable}

\noindent\textbf{Puntuación total sugerida:} 24 puntos $\to$ escalar a la calificación que corresponda.

%======================================================================
\section*{Cierre}
\addcontentsline{toc}{section}{Cierre}
%======================================================================

\begin{cajaTeoria}
\itshape
``El Sistema Solar no es un lugar, es un proceso.''\\
\hfill --- Carl Sagan (paráfrasis)
\end{cajaTeoria}

El simulador es una herramienta para \textbf{hacer preguntas}, no solo para leer respuestas. El catálogo contiene parámetros orbitales referidos a la época J2000 y valores físicos de referencia. Las leyes de Kepler y la gravitación de Newton permiten construir modelos que explican y predicen, dentro de sus supuestos, el movimiento observado. Esta guía te invita a redescubrirlas con tus propias mediciones.

\textbf{Buen viaje orbital.}

%======================================================================
%  BIBLIOGRAFÍA
%======================================================================
\begin{thebibliography}{9}

\bibitem{kepler}
Kepler, J. (1609). \emph{Astronomia Nova}. Praga.

\bibitem{kepler2}
Kepler, J. (1619). \emph{Harmonices Mundi}. Linz.

\bibitem{newton}
Newton, I. (1687). \emph{Philosophiæ Naturalis Principia Mathematica}. Londres.

\bibitem{nasa}
NASA JPL. \emph{Approximate Positions of the Planets}.\\
\url{https://ssd.jpl.nasa.gov/planets/approx_pos.html}

\bibitem{factsheet}
NASA. \emph{Planetary Fact Sheet}.\\
\url{https://nssdc.gsfc.nasa.gov/planetary/factsheet/}

\bibitem{iau}
IAU (2015). \emph{Resolution B3 on Recommended Nominal Conversion Constants for Selected Solar and Planetary Properties}. Honolulu.

\bibitem{vallado}
Vallado, D. A. (2013). \emph{Fundamentals of Astrodynamics and Applications} (4th ed.). Microcosm Press.

\bibitem{murray}
Murray, C. D., \& Dermott, S. F. (1999). \emph{Solar System Dynamics}. Cambridge University Press.

\end{thebibliography}

%======================================================================
%  PÁGINA FINAL
%======================================================================
\vfill
\begin{center}
\rule{0.5\linewidth}{0.5pt}\\[0.3cm]
{\small\textbf{Guía de Ejercicios · Simulador del Sistema Solar}}\\[0.1cm]
{\footnotesize \autor\ · \institucion}\\[0.1cm]
{\footnotesize Distribuido bajo licencia GNU GPL v3}\\[0.1cm]
{\footnotesize \url{https://www.gnu.org/licenses/gpl-3.0}}
\end{center}

\end{document}










